%%
%% Konzept & Umsetzung
%%

\section{Konzept \& Umsetzung}

\subsection{Architektur-Übersicht}

Die Architektur unterstützt zwei Spielmodi.
Im lokalen Modus (Agent vs Human) spielt ein Mensch gegen eine trainierte KI.
Im Multiplayer-Modus (Human vs Human) spielen zwei Menschen über einen Server gegeneinander.
Beide Modi nutzen dieselben GUI-Komponenten, die Spiellogik (Game-Environments und CFR-Solver) stammt aus der Bachelorarbeit und wird unverändert wiederverwendet.

Die Architektur folgt einem Schichtenmodell mit vier Ebenen.
Die \textbf{GUI-Schicht} (PyQt6 in \texttt{src/gui/components/} und \texttt{src/gui/layouts/}) stellt die Darstellung bereit und ist für beide Modi gleich.
Die \textbf{Client-Schicht} gibt es nur im Multiplayer-Modus.
HTTP- und WebSocket-Clients haben ein einheitliches Interface mit denselben Signals, sodass das Protokoll austauschbar ist.
Die \textbf{Server-Schicht} (ebenfalls nur Multiplayer) besteht aus HTTP- und WebSocket-Servern.
Beide bauen auf derselben Klasse \texttt{PokerGameLogic} auf, die Spiellogik ohne transportbezogene Details kapselt und Registrierung, Zustand und Aktionen verwaltet.
Der HTTP-Server nutzt Flask mit REST-Endpoints, der WebSocket-Server läuft asyncio-basiert.
Die \textbf{Logik-Schicht} ist entweder diese geteilte \texttt{PokerGameLogic} auf dem Server oder ein lokales Game-Objekt.

Im lokalen Modus spricht die GUI direkt mit dem Game-Objekt und einem \texttt{StrategyAgent} für die KI.
Es findet keine Netzwerkkommunikation statt.
Im Multiplayer-Modus verbinden sich zwei Clients mit dem Server, erhalten je eine Player-ID (0 oder 1) und tauschen über den Server alle Aktionen und Zustände aus.
Der Server prüft und führt Aktionen aus und verteilt den aktualisierten Spielzustand an beide Clients.

Die Layout-Hierarchie lautet \texttt{BasePokerLayout} $\rightarrow$ \texttt{AgentVsHumanLayout} $\rightarrow$ \texttt{AgentVsHumanGUI} bzw.\ \texttt{HumanVsHumanGUI}.
Konkrete Beispiele für wiederverwendete Komponenten sind \texttt{ActionButtons}, \texttt{PokerTable} und die Player-Widgets.
Entscheidende Prinzipien sind Wiederverwendung (DRY), klare Trennung der Schichten und Thread-Safety im Multiplayer durch \texttt{threading.Lock} in \texttt{PokerGameLogic}.

Die Projektstruktur spiegelt das wider.
In \texttt{src/gui/components/} und \texttt{src/gui/layouts/} liegen die UI-Bausteine, in \texttt{src/gui/server/} Server, Clients und \texttt{game\_logic.py}, in \texttt{src/gui/runner/} die Startskripte.
Verwendet werden PyQt6, Flask, \texttt{websockets} und pytest.

\subsection{Komponenten}

\subsubsection{GUI-Layout}

Die Layout-Hierarchie nutzt Vererbung, um Code-Wiederverwendung zu ermöglichen.
\texttt{BasePokerLayout} stellt die Basis-Funktionalität bereit und erstellt die grundlegende UI-Struktur mit Poker-Tisch, Community-Cards, Pot-Display und Player-Widgets.
\texttt{AgentVsHumanLayout} erweitert \texttt{BasePokerLayout} und fügt spiel-spezifische Komponenten wie Action-Buttons hinzu, sodass beide Spielmodi (\texttt{AgentVsHumanGUI} und \texttt{HumanVsHumanGUI}) die gleiche Basis-Funktionalität nutzen können.

Für die Anordnung der Komponenten verwendet das GUI PyQt6 Layout-Manager: \texttt{QVBoxLayout} ordnet Widgets vertikal an, \texttt{QHBoxLayout} horizontal.
Der Stretch-Faktor wird durch \texttt{addStretch()} eingefügt und erzeugt flexiblen Platz, der sich automatisch an die Fenstergröße anpasst.
Ein kompakter Ausschnitt der Layout-Hierarchie zeigt Listing~\ref{lst:layout}.

\begin{lstlisting}[caption={Layout-Hierarchie (Ausschnitt)}]
class BasePokerLayout(QMainWindow):
    def __init__(self, parent=None):
        super().__init__(parent)
        self.setWindowTitle("CFR Poker Visualizer")
        self.setMinimumSize(1400, 900)
        central_widget = QWidget()
        self.setCentralWidget(central_widget)
        self.main_layout = QVBoxLayout(central_widget)
        self.setup_ui()

    def setup_ui(self):
        self.player_top_widget = TopPlayerWidget(0, "Strategy Agent", self)
        table_container = QWidget()
        self.poker_table = PokerTable()
        # ... Table- und Community-Cards-Setup ...
        self.player_bottom_widget = BottomPlayerWidget(1, "Player", self)
        self.main_layout.addWidget(self.player_top_widget)
        self.main_layout.addWidget(table_container)
        self.main_layout.addWidget(self.player_bottom_widget)

class AgentVsHumanLayout(BasePokerLayout):
    def __init__(self, parent=None):
        super().__init__(parent)
        QTimer.singleShot(100, self.setup_action_buttons)

    def setup_action_buttons(self):
        self.action_buttons = ActionButtons()
        self.control_area.layout().addWidget(self.action_buttons)
\end{lstlisting}\label{lst:layout}

\subsubsection{WebSocket Client-Server-Kommunikation}

Die WebSocket-Kommunikation nutzt asyncio für Push-basierte Updates statt Polling.
Der Server sendet State-Updates automatisch an alle verbundenen Clients, sobald sich der Spielzustand ändert.
Der Client verwendet einen separaten \texttt{QThread} mit eigener asyncio-Event-Loop, um die Qt-Event-Loop nicht zu blockieren.
Auf dem Server genügt nach Registrierung eines Clients der Versand der \texttt{player\_id} und bei jeder Zustandsänderung ein Broadcast an alle Verbindungen; eingehende Aktionen werden in \texttt{handle\_action} verarbeitet und lösen erneut einen State-Broadcast aus.
Im Client (Listing~\ref{lst:ws-client}) laufen asyncio und Qt getrennt: Der Worker-Thread startet eine eigene Event-Loop, empfängt Nachrichten und leitet sie per PyQt-Signal an die GUI weiter.

\begin{lstlisting}[caption={WebSocket-Client: QThread mit asyncio-Event-Loop}]
class WebSocketWorker(QThread):
    state_received = pyqtSignal(dict)
    connected_signal = pyqtSignal(int)

    def run(self):
        self.loop = asyncio.new_event_loop()
        asyncio.set_event_loop(self.loop)
        try:
            self.loop.run_until_complete(self._connect_and_listen())
        finally:
            self.loop.close()

    async def _connect_and_listen(self):
        self.websocket = await connect(self.server_url)
        msg = await self.websocket.recv()
        data = json.loads(msg)
        self.player_id = data.get("player_id")
        self.connected_signal.emit(self.player_id)
        async for message in self.websocket:
            data = json.loads(message)
            if data.get("type") == "state":
                self.state_received.emit(data["state"])
\end{lstlisting}\label{lst:ws-client}

Die Integration von asyncio mit Qt erfolgt durch diesen separaten \texttt{QThread}, der eine eigene asyncio-Event-Loop besitzt.
Async-Events werden über PyQt6 Signals an den GUI-Thread übertragen, wodurch eine saubere Trennung zwischen asynchroner Netzwerk-Kommunikation und GUI-Updates erreicht wird.
Dies ermöglicht Push-basierte Kommunikation ohne Blockierung der GUI und verbessert die Effizienz erheblich im Vergleich zu HTTP-Polling.

\subsubsection{Tests}

Die Tests folgen einem szenarien-basierten Ansatz, der typische Anwendungsfälle abdeckt.
Dies macht die Tests verständlicher und aussagekräftiger als viele kleine Edge-Case-Tests.
Für WebSocket-Tests wird \textbf{pytest-asyncio} verwendet, um asynchrone Test-Funktionen zu ermöglichen (Listing~\ref{lst:test}).

\begin{lstlisting}[caption={Szenarien-basierter WebSocket-Test (Ausschnitt)}]
@pytest.fixture(scope="module")
def ws_server_url():
    port = _free_port()
    game = KuhnPokerGame()
    server = PokerWebSocketServer(game, host="localhost", port=port)
    # ... Server im Thread starten ...
    yield f"ws://localhost:{port}"

async def test_zwei_clients_verbinden_bekommen_player_id_und_state(ws_server_url):
    async with websockets.connect(ws_server_url) as ws1, \
                 websockets.connect(ws_server_url) as ws2:
        msgs1 = await _recv_n(ws1, 3)
        msgs2 = await _recv_n(ws2, 3)
    assert next(m for m in msgs1 if m.get("type") == "player_id")["player_id"] == 0
    assert next(m for m in msgs2 if m.get("type") == "player_id")["player_id"] == 1
    state_msg = next(m for m in msgs1 if m.get("type") == "state")
    assert "legal_actions" in state_msg["state"] and "current_player" in state_msg["state"]
\end{lstlisting}\label{lst:test}

Die Tests decken die wichtigsten Szenarien ab: Client-Registrierung, State-Updates, Action-Verarbeitung und Fehlerbehandlung.
GUI-Darstellung und Animationen wurden manuell geprüft, da diese schwer automatisiert werden können.
\newpage
